\documentclass[11pt]{article}
\usepackage{explorer}
\newcommand{\ExplorerLanguage}{FR}
\newcommand{\ExplorerNumber}{10}
\newcommand{\ExplorerKind}{Leçon et méthode}
\begin{document}
\ExplorerTitle{Première--Terminale · suites}{Sommes télescopiques}
\ExplorerSection{Avant de commencer}
Fractions, sommes finies et indices.\par
\ExplorerSection{Idées et méthode}
Une somme télescopique s’écrit comme une somme de différences consécutives. Les termes intérieurs s’annulent : \(\sum (a_{k}-a_{k+1})=a_{1}-a_{n+1}\) pour k allant de 1 à n.\par
Écrire d’abord trois termes et les deux derniers permet de voir les annulations et de contrôler les bornes. Une simplification terme à terme ne dispense pas de préciser les indices.\par
Une décomposition utile est \(\frac{1}{k(k+1)}=\frac{1}{k}-\frac{1}{k+1}\), valable pour \(k\geq 1\). On la vérifie en mettant les fractions au même dénominateur.\par
\ExplorerSection{Exemple expliqué}
\(\sum  \frac{1}{k(k+1)}\), de \(k=1\) à n, vaut \(1-\frac{1}{n+1}=\frac{n}{n+1}\). La somme est inférieure à 1 et tend vers 1 lorsque n augmente.\par
\ExplorerSection{À explorer}
Calculer les cinq premières sommes partielles, puis étudier l’écart à 1. Distinguer la somme finie de sa limite.\par
\end{document}
